% Options for packages loaded elsewhere
\PassOptionsToPackage{unicode}{hyperref}
\PassOptionsToPackage{hyphens}{url}
%
\documentclass[
  10pt,
  letterpaper,
]{article}
\usepackage{amsmath,amssymb}
\usepackage{iftex}
\ifPDFTeX
  \usepackage[T1]{fontenc}
  \usepackage[utf8]{inputenc}
  \usepackage{textcomp} % provide euro and other symbols
\else % if luatex or xetex
  \usepackage{unicode-math} % this also loads fontspec
  \defaultfontfeatures{Scale=MatchLowercase}
  \defaultfontfeatures[\rmfamily]{Ligatures=TeX,Scale=1}
\fi
\usepackage{lmodern}
\ifPDFTeX\else
  % xetex/luatex font selection
\fi
% Use upquote if available, for straight quotes in verbatim environments
\IfFileExists{upquote.sty}{\usepackage{upquote}}{}
\IfFileExists{microtype.sty}{% use microtype if available
  \usepackage[]{microtype}
  \UseMicrotypeSet[protrusion]{basicmath} % disable protrusion for tt fonts
}{}
\makeatletter
\@ifundefined{KOMAClassName}{% if non-KOMA class
  \IfFileExists{parskip.sty}{%
    \usepackage{parskip}
  }{% else
    \setlength{\parindent}{0pt}
    \setlength{\parskip}{6pt plus 2pt minus 1pt}}
}{% if KOMA class
  \KOMAoptions{parskip=half}}
\makeatother
\usepackage{xcolor}
\usepackage[margin=1in]{geometry}
\usepackage{color}
\usepackage{fancyvrb}
\newcommand{\VerbBar}{|}
\newcommand{\VERB}{\Verb[commandchars=\\\{\}]}
\DefineVerbatimEnvironment{Highlighting}{Verbatim}{commandchars=\\\{\}}
% Add ',fontsize=\small' for more characters per line
\newenvironment{Shaded}{}{}
\newcommand{\AlertTok}[1]{\textcolor[rgb]{1.00,0.00,0.00}{\textbf{#1}}}
\newcommand{\AnnotationTok}[1]{\textcolor[rgb]{0.38,0.63,0.69}{\textbf{\textit{#1}}}}
\newcommand{\AttributeTok}[1]{\textcolor[rgb]{0.49,0.56,0.16}{#1}}
\newcommand{\BaseNTok}[1]{\textcolor[rgb]{0.25,0.63,0.44}{#1}}
\newcommand{\BuiltInTok}[1]{\textcolor[rgb]{0.00,0.50,0.00}{#1}}
\newcommand{\CharTok}[1]{\textcolor[rgb]{0.25,0.44,0.63}{#1}}
\newcommand{\CommentTok}[1]{\textcolor[rgb]{0.38,0.63,0.69}{\textit{#1}}}
\newcommand{\CommentVarTok}[1]{\textcolor[rgb]{0.38,0.63,0.69}{\textbf{\textit{#1}}}}
\newcommand{\ConstantTok}[1]{\textcolor[rgb]{0.53,0.00,0.00}{#1}}
\newcommand{\ControlFlowTok}[1]{\textcolor[rgb]{0.00,0.44,0.13}{\textbf{#1}}}
\newcommand{\DataTypeTok}[1]{\textcolor[rgb]{0.56,0.13,0.00}{#1}}
\newcommand{\DecValTok}[1]{\textcolor[rgb]{0.25,0.63,0.44}{#1}}
\newcommand{\DocumentationTok}[1]{\textcolor[rgb]{0.73,0.13,0.13}{\textit{#1}}}
\newcommand{\ErrorTok}[1]{\textcolor[rgb]{1.00,0.00,0.00}{\textbf{#1}}}
\newcommand{\ExtensionTok}[1]{#1}
\newcommand{\FloatTok}[1]{\textcolor[rgb]{0.25,0.63,0.44}{#1}}
\newcommand{\FunctionTok}[1]{\textcolor[rgb]{0.02,0.16,0.49}{#1}}
\newcommand{\ImportTok}[1]{\textcolor[rgb]{0.00,0.50,0.00}{\textbf{#1}}}
\newcommand{\InformationTok}[1]{\textcolor[rgb]{0.38,0.63,0.69}{\textbf{\textit{#1}}}}
\newcommand{\KeywordTok}[1]{\textcolor[rgb]{0.00,0.44,0.13}{\textbf{#1}}}
\newcommand{\NormalTok}[1]{#1}
\newcommand{\OperatorTok}[1]{\textcolor[rgb]{0.40,0.40,0.40}{#1}}
\newcommand{\OtherTok}[1]{\textcolor[rgb]{0.00,0.44,0.13}{#1}}
\newcommand{\PreprocessorTok}[1]{\textcolor[rgb]{0.74,0.48,0.00}{#1}}
\newcommand{\RegionMarkerTok}[1]{#1}
\newcommand{\SpecialCharTok}[1]{\textcolor[rgb]{0.25,0.44,0.63}{#1}}
\newcommand{\SpecialStringTok}[1]{\textcolor[rgb]{0.73,0.40,0.53}{#1}}
\newcommand{\StringTok}[1]{\textcolor[rgb]{0.25,0.44,0.63}{#1}}
\newcommand{\VariableTok}[1]{\textcolor[rgb]{0.10,0.09,0.49}{#1}}
\newcommand{\VerbatimStringTok}[1]{\textcolor[rgb]{0.25,0.44,0.63}{#1}}
\newcommand{\WarningTok}[1]{\textcolor[rgb]{0.38,0.63,0.69}{\textbf{\textit{#1}}}}
\setlength{\emergencystretch}{3em} % prevent overfull lines
\providecommand{\tightlist}{%
  \setlength{\itemsep}{0pt}\setlength{\parskip}{0pt}}
\setcounter{secnumdepth}{-\maxdimen} % remove section numbering
\usepackage{xurl}
\setlength{\emergencystretch}{3em}
\usepackage{microtype}
\ifLuaTeX
  \usepackage{selnolig}  % disable illegal ligatures
\fi
\usepackage{bookmark}
\IfFileExists{xurl.sty}{\usepackage{xurl}}{} % add URL line breaks if available
\urlstyle{same}
\hypersetup{
  hidelinks,
  pdfcreator={LaTeX via pandoc}}

\author{}
\date{}

\begin{document}

\section{DDTA BAProposition structural candidate -
R1}\label{ddta-baproposition-structural-candidate---r1}

\textbf{Status:} CANDIDATE / NOT ACCEPTED / BA2 OPEN

\textbf{Derived by:} BA2-T1

\textbf{Repository baseline reviewed:}
\texttt{3d8251328c77177375cccf1c51caa54b7473e21e}

\textbf{Closed BA1 dependency:} \texttt{BAReferent\ +\ BAProposition}

\subsection{1. Candidate purpose}\label{candidate-purpose}

This artifact proposes only the minimum internal semantic shape that a
\texttt{BAProposition} appears to need after BA1 closure. It does not
close the canonical relation/action vocabulary.

The lower-bound candidate is:

\begin{Shaded}
\begin{Highlighting}[]
\NormalTok{BAProposition                         [identity accepted by BA1]}
\NormalTok{|{-} semanticOperator                  1  [explicit method{-}neutral semantic key]}
\NormalTok{|{-} participation                     1..*}
\NormalTok{|    |{-} role                         1  [explicit participation{-}role key]}
\NormalTok{|    \textasciigrave{}{-} term                         1  [BAReferent, or typed local value when identity is not required]}
\NormalTok{\textasciigrave{}{-} scopedModifier                    0..* [required capability; exact material form still open]}
\end{Highlighting}
\end{Shaded}

\texttt{semanticOperator}, \texttt{role} and \texttt{scopedModifier} are
\textbf{not new BAE identity families}. Their canonical vocabularies and
exact encoding remain BA2 work.

\subsection{2. Why pure binary SPO is not the universal
core}\label{why-pure-binary-spo-is-not-the-universal-core}

The documentation layer already states that SPO references expose
reusable participants but do not contain all conditional, concurrency,
lifecycle or failure semantics.

The facial-access \texttt{FR-3.4} meaning involves, in one reusable
delivery fact:

\begin{itemize}
\tightlist
\item
  source/responsible endpoint;
\item
  destination;
\item
  delivered information;
\item
  request correlation;
\item
  completion/failure semantics.
\end{itemize}

The order-fulfillment corpus adds all-or-nothing reservation,
concurrency, lifecycle, compensation and indeterminate result semantics.

A universal \texttt{subject-predicate-object} triple would therefore
require either:

\begin{enumerate}
\def\labelenumi{\arabic{enumi}.}
\tightlist
\item
  fragmenting one atomic claim into several relations and later
  reconstructing their common scope; or
\item
  inventing a synthetic referent for every relation occurrence merely to
  hang the missing roles/qualifiers on it.
\end{enumerate}

The first loses explicit assertion scope; the second conflicts with the
BA1 rule that a \texttt{BAReferent} denotes independently reusable
project meaning rather than serving as a generic proxy for an assertion
record.

Therefore binary relations remain a valid \textbf{special
case/projection}, but not the only admissible BAProposition shape.

\subsection{\texorpdfstring{3. \texttt{semanticOperator} lower
bound}{3. semanticOperator lower bound}}\label{semanticoperator-lower-bound}

A proposition requires an explicit semantic operator/predicate key so
that a consumer does not need to reconstruct from prose what kind of
assertion is being made.

Examples such as \texttt{delivers}, \texttt{correlates},
\texttt{dependsOn}, \texttt{realizes}, \texttt{constrains},
\texttt{transitions} or \texttt{produces} are illustrative only. BA2-T1
does \textbf{not} accept these labels as the final vocabulary.

The lower-bound invariant is only:

\begin{quote}
every BAProposition must expose a stable method-neutral operator
identity/key sufficient to distinguish assertion semantics mechanically.
\end{quote}

The operator key is vocabulary identity, not a third BAE family.

\subsection{4. Explicit role-bound
participation}\label{explicit-role-bound-participation}

Each proposition must expose its participants through explicit role
bindings rather than anonymous positional ordering.

Conceptually:

\begin{Shaded}
\begin{Highlighting}[]
\NormalTok{ParticipationBinding}
\NormalTok{|{-} role}
\NormalTok{\textasciigrave{}{-} term}
\end{Highlighting}
\end{Shaded}

The \texttt{term} is normally a \texttt{BAReferent}. A typed local value
may be used when the semantic content does not need independent
cross-proposition identity; if it does, BA1 requires promotion of the
project meaning to \texttt{BAReferent} identity.

The same BAReferent may:

\begin{itemize}
\tightlist
\item
  participate in multiple propositions;
\item
  occupy different roles in different propositions;
\item
  occur more than once when the operator semantics genuinely allow it.
\end{itemize}

Role cardinality and operator-role compatibility are not closed by T1.

\subsection{5. Why roles are not semantic
kinds}\label{why-roles-are-not-semantic-kinds}

Participation role is contextual. Referent classification is
intrinsic/shared enough to be reusable across propositions.

For example, \texttt{RecognitionCapture} may be classified
method-neutrally as information/resource-like project meaning, while
participating as transferred information in one proposition and as an
accepted/validated object in another.

Therefore:

\begin{Shaded}
\begin{Highlighting}[]
\NormalTok{referent semantic classification}
\NormalTok{        !=}
\NormalTok{proposition participation role}
\end{Highlighting}
\end{Shaded}

BA2-T1 supports an explicit method-neutral classification capability for
mechanical projection, but does not freeze the classification
vocabulary, cardinality or storage location.

A consumer must not have to infer \texttt{information},
\texttt{capability}, \texttt{behavior}, \texttt{contract},
\texttt{state} or similar reusable distinctions again from raw prose
merely because roles happen to suggest them.

\subsection{6. Action/behavior identity
boundary}\label{actionbehavior-identity-boundary}

BA1 already allows behavior/event meaning to be a \texttt{BAReferent}
when independently reusable.

BA2-T1 therefore keeps two concerns separate:

\begin{itemize}
\tightlist
\item
  a concrete/reusable behavior or event may have \texttt{BAReferent}
  identity;
\item
  a BAProposition's \texttt{semanticOperator} expresses what the current
  assertion says about its participants.
\end{itemize}

No new \texttt{Action}, \texttt{Event}, \texttt{Relation} or
\texttt{Participation} BAE family is forced.

\subsection{7. Scoped modifier
capability}\label{scoped-modifier-capability}

Conditions, states, failure rules, applicability, atomicity, concurrency
and realization semantics can change the truth/meaning of a proposition
without necessarily introducing another independently reusable project
referent.

BA2-T1 therefore requires a \textbf{non-ambiguous proposition-scoped
modifier capability}.

It does not yet decide the final material form. A later BA2 step must
pressure-test whether a modifier is best represented as:

\begin{itemize}
\tightlist
\item
  a role-bound referent/value in the same n-ary proposition;
\item
  a normalized embedded qualifier record;
\item
  a separate BAProposition about a reusable context/referent;
\item
  or a constrained combination of these.
\end{itemize}

What is rejected is an untyped free-text ``extra semantics'' bag that
forces consumers back to prose interpretation.

\subsection{8. Illustrative facial-access
shape}\label{illustrative-facial-access-shape}

The labels below are illustrative and \textbf{not accepted canonical
vocabulary}:

\begin{Shaded}
\begin{Highlighting}[]
\NormalTok{P{-}FR34{-}DELIVERY}
\NormalTok{operator: transfer}
\NormalTok{participants:}
\NormalTok{  {-} role: source              {-}\textgreater{} CameraSubsystem}
\NormalTok{  {-} role: destination         {-}\textgreater{} RecognitionProcessor}
\NormalTok{  {-} role: information         {-}\textgreater{} RecognitionCapture}
\NormalTok{  {-} role: correlationContext  {-}\textgreater{} RecognitionRequest}
\NormalTok{modifiers:}
\NormalTok{  {-} completion: incomplete delivery != successful completion}
\end{Highlighting}
\end{Shaded}

If the delivery behavior itself needs reuse by security constraints,
that project meaning also receives a \texttt{BAReferent} identity; the
proposition structure does not use its own assertion identity as a proxy
for that behavior.

\subsection{9. Illustrative order-fulfillment
pressure}\label{illustrative-order-fulfillment-pressure}

A full-order reservation clause cannot be faithfully reduced to one
anonymous triple when it simultaneously connects an executor, request,
evaluation correlation, reservation/result meaning and
atomic/concurrency conditions.

A lower-bound proposition can bind those participants explicitly while
keeping atomicity/concurrency as scoped semantics. The exact canonical
operator and role names remain open.

\subsection{10. Alternatives rejected or
deferred}\label{alternatives-rejected-or-deferred}

\begin{itemize}
\tightlist
\item
  \textbf{Pure SPO as universal shape:} REJECTED.
\item
  \textbf{Binary propositions allowed as special cases:} SUPPORTED.
\item
  \textbf{Anonymous positional participant list:} REJECTED.
\item
  \textbf{Explicit role-bound n-ary participation:} FORCED CANDIDATE.
\item
  \textbf{Fixed DFD-style roles (\texttt{process}, \texttt{dataflow},
  \texttt{store}, etc.):} REJECTED AS UNIVERSAL CORE.
\item
  \textbf{Stable method-neutral semantic operator key:} FORCED
  CANDIDATE.
\item
  \textbf{Exhaustive operator vocabulary:} DEFERRED.
\item
  \textbf{Participation/Role/Qualifier as new BA1 identity families:}
  NOT FORCED.
\item
  \textbf{Scoped modifier capability:} REQUIRED; exact encoding OPEN.
\item
  \textbf{Referent classification inferred only from proposition roles:}
  REJECTED.
\item
  \textbf{Explicit reusable method-neutral classification capability:}
  SUPPORTED / VOCABULARY OPEN.
\item
  \textbf{BA3 provenance fields:} DEFERRED TO BA3.
\end{itemize}

\subsection{11. Falsification rule}\label{falsification-rule}

This candidate must be revised if a concrete governed corpus or
method-neutral consumer demonstrates that:

\begin{enumerate}
\def\labelenumi{\arabic{enumi}.}
\tightlist
\item
  explicit operator + role-bound n-ary participation cannot preserve an
  atomic project fact without hidden prose reconstruction;
\item
  proposition-scoped modifiers cannot be represented without creating an
  unacknowledged third identity family;
\item
  referent classification cannot remain separate from participation
  roles;
\item
  a fixed binary-only structure preserves all reviewed semantics more
  minimally without synthetic referent/proposition grouping; or
\item
  a required role/operator distinction is actually method-specific
  rather than methodology-neutral.
\end{enumerate}

\subsection{12. Next step}\label{next-step}

\textbf{BA2-T2 - semantic operator, participation-role and
scoped-modifier vocabulary pressure test.}

BA2 remains open. BA3 remains not started.

\end{document}
